\section{Methods}

Both studies use the same core invasion protocol. A population of strategies is evaluated over generations, weaker strategies are removed, and new strategies are injected by a specified generator. Selection pressure is applied in train conditions and robustness is reported in held-out test conditions. This separates adaptation from evaluation and turns governance comparisons into tests of out-of-distribution stability rather than only in-sample fit.

\subsection{Study 1: Fishery Commons as a signal-design precursor}

Fishery Commons is a minimal renewable commons with one shared stock, per-agent harvest actions, regeneration, and collapse dynamics. Governance changes extraction pressure through monitoring, sanctions, quota clipping, adaptive quotas, or temporary closures. The first Fishery layer is a matched baseline that compares no governance against monitoring with sanctions under two injector families: mutation and live JSON-based LLM injection. The second layer is a broader non-LLM top-down comparison in the medium tier that evaluates monitoring with sanctions, adaptive quotas, and temporary closures across balanced and adversarial-heavy partner mixes, adversarial pressures \(0.3\) and \(0.5\), and three injector families: mutation, an adversarial heuristic injector, and search over mutations. A medium-tier PPO self-play validation is retained as a policy-family robustness check on the resulting ranking.

\subsection{Study 2: Harvest Commons as the main architecture study}

Harvest Commons extends the commons problem from one global stock to local renewable patches with local interaction, optional communication, and optional side-payments. Agents act through parameterized local harvest, restraint, request, offer, and compliance components, while governance is imposed by the environment. Across the broader project history, the architecture space includes \texttt{none}, \texttt{bottom\_up\_only}, \texttt{top\_down\_only}, and \texttt{hybrid}. In this follow-up, those conditions are interpreted as institutional designs rather than as simple environment toggles.

The architecture evidence is organized in two stages. The first is an architecture-restoration pilot over \texttt{none}, \texttt{bottom\_up\_only}, \texttt{top\_down\_only}, and \texttt{hybrid} in the medium and hard tiers, balanced and adversarial-heavy partner mixes, pressure \(0.3\), and two non-LLM attackers: mutation and search over mutations. The second is the high-power confirmatory comparison over \texttt{bottom\_up\_only}, \texttt{top\_down\_only}, and \texttt{hybrid} under the stronger \texttt{search\_mutation} attacker in the medium and hard tiers, balanced and adversarial-heavy partner mixes, and pressures \(0.3\) and \(0.5\).

The confirmatory ranking rule is lexicographic: lower garden failure first, then higher patch health, then higher welfare. Accordingly, a condition ``wins'' a decision cell only in the sense of first rank under that rule. Pairwise contrasts are reported for \texttt{hybrid - top\_down\_only}, \texttt{hybrid - bottom\_up\_only}, and \texttt{top\_down\_only - bottom\_up\_only}.

\subsection{Welfare incidence and capability ladder}

Harvest evaluation logs per-agent realized welfare, requested harvest, realized harvest, prevented harvest, credit sent, credit received, local patch health, and whether the agent was targeted or capped. These logs support two incidence views. The first groups agents into low-, middle-, and high-aggression terciles using aggressive-request behavior, with requested harvest as a fallback grouping variable. The second groups agents by whether they were targeted for caps during the episode. These summaries are used to distinguish concentrated restraint from broad welfare taxation.

Attacker strength is represented by a non-LLM capability ladder that progresses from random injection to mutation, then to an adversarial heuristic attacker, and finally to search over mutations. The ladder pilot evaluates \texttt{bottom\_up\_only}, \texttt{top\_down\_only}, and \texttt{hybrid} in the medium and hard tiers, balanced and adversarial-heavy partner mixes, pressures \(0.3\) and \(0.5\), and three runs per cell. In the follow-up, these attackers are interpreted as increasing rungs of adversarial capability. The ladder analysis records the first rung at which a pairwise architecture contrast loses its patch-health advantage, the first rung at which it loses its neighborhood-overharvest advantage, and the first rung at which the welfare cost crosses the \(-0.15\) threshold.
